% curve25515_impl.tex — replacement for \subsection{Implementation} of
% \section{Case Study 2: Curve25519 in Microcode} (label sec:curve25519) in
% gold489.tex. Drop in with % curve25515_impl.tex — replacement for \subsection{Implementation} of
% \section{Case Study 2: Curve25519 in Microcode} (label sec:curve25519) in
% gold489.tex. Drop in with % curve25515_impl.tex — replacement for \subsection{Implementation} of
% \section{Case Study 2: Curve25519 in Microcode} (label sec:curve25519) in
% gold489.tex. Drop in with % curve25515_impl.tex — replacement for \subsection{Implementation} of
% \section{Case Study 2: Curve25519 in Microcode} (label sec:curve25519) in
% gold489.tex. Drop in with \input{curve25515_impl} in place of the existing
% subsection, or paste inline.
%
% Every triad, uop and slot-density figure below is taken from the installed
% patch arrays in simple/curve25519/full_curve25519_inline2.c.
%
% NOTE: the original text cited \Cref{sec:ucode-layer}, which is not defined
% anywhere in gold489.tex. This version states the intra-triad forwarding
% property locally and cites \Cref{sec:cyclewalk}, which establishes the same
% slot semantics for Keccak. When \section{Background and Programming Model}
% (sec:background) is written, that claim probably belongs there, cited from
% both case studies.

\subsection{Implementation}

\subsubsection{Register interface and operand lifetimes}
\label{sec:c25519-interface}

Neither patch performs a memory access.
All operands, accumulators, carries, and results live in registers for the whole
invocation, and the surrounding x86 wrapper is responsible for placing the inputs
and retrieving the outputs.
This division follows from the resource that binds each side.
Patch RAM is scarce and native instruction slots are not, so we move every
operation that does not require the microcode interface into the wrapper.
\Cref{tab:c25519-regif} gives the resulting register interface.

The wrapper also supplies two pre-zeroed accumulator seeds in \code{RAX} and
\code{R8}, which removes two initialisation $\mu$ops from each patch.
For multiplication the wrapper issues ten loads and two register clears.
For squaring it issues five loads, four \code{LEA} instructions that form the
doubled operands, two \code{IMUL} instructions that form the operands scaled by
19, and the same two clears.

\begin{table}[htbp]
\centering
\caption{Register interface of the two field-arithmetic patches.
The first operand of multiplication occupies five registers that no
$\mu$op overwrites, so all 25 partial products read them in place.
The second operand is copied to \code{TMP10}--\code{TMP14} at patch entry, which
frees its five architectural registers to receive output limbs as successive
positions complete.
The squaring patch receives the doubled and scaled operands from the wrapper and
returns its result in a different register set, which \Cref{sec:c25519-integration}
uses to explain where chaining is possible.}
\label{tab:c25519-regif}
\footnotesize
\setlength{\tabcolsep}{4pt}
\begin{tabularx}{\linewidth}{@{}llX@{}}
\toprule
Role & Registers & Lifetime \\
\midrule
\multicolumn{3}{@{}l}{\emph{multiplication}} \\
$a_0,\ldots,a_4$      & \code{RDI RSI R12 R11 R14} & read-only for the entire patch \\
$b_0,\ldots,b_4$      & \code{R15 R13 R9 R10 RBX}  & copied to \code{TMP10}--\code{TMP14}, then reused as outputs \\
$19b_1,\ldots,19b_4$  & \code{R13 R9 R10 RBX}      & formed in place at entry, consumed in decreasing index \\
$h_0,\ldots,h_4$      & \code{R15 R13 R9 R10 RAX}  & written as each output position completes \\
\midrule
\multicolumn{3}{@{}l}{\emph{squaring}} \\
$a_0,\ldots,a_4$      & \code{RDI RSI R12 R11 R14} & \code{RDI} is overwritten by $h_0$ at the end of position zero \\
$2a_0,\ldots,2a_3$    & \code{R15 R13 R9 R10}      & supplied by the wrapper; \code{R13} is regenerated inside the patch \\
$19a_4$, $19a_3$      & \code{RBX RDX}             & supplied by the wrapper \\
$h_0,\ldots,h_4$      & \code{RDI R9 R10 RBX RAX}  & written as each output position completes \\
\midrule
\multicolumn{3}{@{}l}{\emph{shared internal state}} \\
low accumulator       & \code{TMP0}   & holds the incoming carry, then the sum of the low product halves \\
high accumulator      & \code{R8}     & sum of the high product halves \\
carry accumulator     & \code{TMP9}   & count of carries out of the low accumulator \\
materialised carry    & \code{TMP15}  & destination of \code{SETCC\_CONDB} \\
product high half     & \code{RCX}    & second output of every multiply \\
product staging       & \code{RDX}    & the operand each multiply is allowed to destroy \\
limb extraction       & \code{TMP1 TMP2 TMP8} & shift results combined into the next carry \\
\bottomrule
\end{tabularx}
\end{table}

The unusual feature of this interface is the treatment of the second operand,
and it follows from the shape of the multiply $\mu$op.
\code{MUL\_DSZ64\_DRR} has two destinations.
It preserves its first source, writes the low half of the product over its
\emph{second} source, and writes the high half to the named destination.
Every product therefore destroys one of the operands that formed it.
A schoolbook multiplication reads each limb of both operands five times, so no
limb can be left in the register that a product will consume.

We resolve this by separating an immortal copy of each operand from a mortal
staging register.
The first operand is immortal by construction, since it always occupies the
preserved source position.
The second operand is copied to \code{TMP10}--\code{TMP14} in the first triad of
the patch, and thereafter each product is formed by copying the required limb into
\code{RDX} and letting the multiply overwrite it.
The three-triad prologue performs this copy and forms $19b_1,\ldots,19b_4$ in
place in the architectural registers, using the same destructive low-half write
that the schoolbook products use.
Both forms of the second operand are consequently live at once, the unscaled form
in temporary registers and the scaled form in architectural ones.

\subsubsection{Register-resident multiplication}
\label{sec:c25519-mul}

The multiplication patch occupies 66 triads and contains 196 $\mu$ops, an average
of 2.97 of the three available slots.
It is straight-line code with no microsequencer control flow.
This distinguishes it from the Keccak kernel of \Cref{sec:keccak}, where a loop
was unavoidable because a full permutation cannot be represented in 128 triads.
A field operation contains no repetition worth keeping resident, and its complete
body fits, so we spend no triad on loop control and no $\mu$op on a loop variable.

The patch computes the five output positions in sequence.
Three triads form the prologue, twelve triads compute each output position, and
three triads perform the final reduction.
For position $k$ we accumulate the five products required by the recurrence of
\Cref{sec:c25519-strategy}, extract the low 51 bits as limb $k$, and pass the
remaining bits to position $k+1$.

\paragraph{The product staircase.}
The five twelve-triad blocks are the same schedule with different operands.
This regularity is not incidental, and \Cref{tab:c25519-staircase} shows where it
comes from.
The first operand advances through $a_0,\ldots,a_4$ in the same order in every
block, so its five registers are read in a fixed sequence.
Only the second operand sequence moves, and it moves by exactly one position per
block: it begins at $b_k$, descends through the unscaled limbs to $b_0$, and then
continues into the scaled limbs from $19b_4$ downwards.
Position four requires no scaled term, because $i+j=9$ has no solution with
$i,j\leq4$, and position three requires only the single term $19a_4b_4$.

\begin{table}[htbp]
\centering
\caption{Second-operand sequence for each output position of the multiplication
patch. The first operand is always read as $a_0,a_1,a_2,a_3,a_4$ in that order,
so each row lists only the register staged into \code{RDX} for the corresponding
product. Unscaled limbs are read from the temporary copies and scaled limbs from
the architectural registers. The rightmost column names the register that
receives the completed output limb, which in every case is the architectural
register whose scaled multiple was last required by the preceding row.}
\label{tab:c25519-staircase}
\footnotesize
\setlength{\tabcolsep}{5pt}
\begin{tabular}{@{}clllllc@{}}
\toprule
$k$ & $\times a_0$ & $\times a_1$ & $\times a_2$ & $\times a_3$ & $\times a_4$ & $h_k$ \\
\midrule
0 & $b_0$ & $19b_4$ & $19b_3$ & $19b_2$ & $19b_1$ & \code{R15} \\
1 & $b_1$ & $b_0$    & $19b_4$ & $19b_3$ & $19b_2$ & \code{R13} \\
2 & $b_2$ & $b_1$    & $b_0$    & $19b_4$ & $19b_3$ & \code{R9}  \\
3 & $b_3$ & $b_2$    & $b_1$    & $b_0$    & $19b_4$ & \code{R10} \\
4 & $b_4$ & $b_3$    & $b_2$    & $b_1$    & $b_0$    & \code{RAX} \\
\bottomrule
\end{tabular}
\end{table}

The staircase determines where the output limbs are placed.
The scaled multiple $19b_j$ is required for the last time in the block for
position $4-j$, and the register holding it is dead from that point onward.
We write the output limb of that block into precisely that register.
Limb zero therefore lands in \code{R15}, which held $b_0$ and whose value is
preserved in \code{TMP10}, and limbs one through three land in \code{R13},
\code{R9} and \code{R10} as $19b_1$, $19b_2$ and $19b_3$ retire in turn.
Only limb four requires a register that never held part of an operand, and it
uses \code{RAX}, which the wrapper supplies as a zeroed seed and which is free once
the low accumulator has been initialised.
The patch thus completes a full multiplication without ever holding more than one
copy of any operand limb and without a single spill.

\paragraph{Three accumulators and materialised carries.}
Each output position sums five 128-bit products.
We represent the running sum in three registers rather than two.
\code{TMP0} accumulates the low halves, \code{R8} accumulates the high halves, and
\code{TMP9} counts the carries that leave \code{TMP0}.

The third accumulator is what the microcode carry interface buys us.
On Goldmont, each arithmetic $\mu$op records its carry with its destination
register rather than in a single architectural flag, and
\code{SETCC\_CONDB\_DR} converts that recorded carry into a zero or one value.
We place the \code{SETCC} in the slot immediately after the addition that produces
the carry, so the two occupy one triad and the materialisation costs no additional
triad.
The materialised value is then added to \code{TMP9} while the high-half chain
proceeds in a different slot.
A conventional \code{ADC} chain would instead serialise the two halves, because
\code{ADC} reads the architectural carry flag and only one carry can be pending in
that flag at a time.
Goldmont implements neither ADX nor a second carry flag, so native code for this
representation has no equivalent of the third accumulator.
Keeping the carry count separate costs one triad per output position and removes a
serialisation point from the dependency path of all five.

Two properties of the microcode interface constrain how this chain is written.
\code{SETCC\_CONDB\_DR} writes a correct value only when its destination is a
temporary register, so all three accumulators and the materialised carry are
temporary registers and the chain never touches architectural state.
The recorded carry also persists only until the next arithmetic operation writes
the same destination, which fixes the position of the \code{SETCC} relative to its
addition and is the reason the two are packed together rather than scheduled
apart.

\begin{lstlisting}[caption={The opening five triads of an output position, taken
from position one. Slot~2 of the first triad forms the incoming carry, slot~0 of
each even triad stages the next second-operand limb into \code{RDX}, and each
addition into \code{TMP0} is immediately followed by the \code{SETCC} that
materialises its carry. The high accumulator and the carry accumulator are
initialised in the third triad, after which both chains advance in parallel with
the multiplies.},label={lst:c25519-mulblock}]
{ ZEROEXT_DSZ64_DR(RDX, TMP11), MUL_DSZ64_DRR(RCX, RDI, RDX),
  OR_DSZ64_DRR(TMP0, TMP8, TMP1), NOP_SEQWORD },
{ ADD_DSZ64_DRR(TMP0, TMP0, RDX), SETCC_CONDB_DR(TMP15, TMP0),
  ZEROEXT_DSZ64_DR(RDX, TMP10), NOP_SEQWORD },
{ ZEROEXT_DSZ64_DR(R8, RCX), MUL_DSZ64_DRR(RCX, RSI, RDX),
  ZEROEXT_DSZ64_DR(TMP9, TMP15), NOP_SEQWORD },
{ ADD_DSZ64_DRR(TMP0, TMP0, RDX), SETCC_CONDB_DR(TMP15, TMP0),
  ZEROEXT_DSZ64_DR(RDX, RBX), NOP_SEQWORD },
{ ADD_DSZ64_DRR(R8, R8, RCX), MUL_DSZ64_DRR(RCX, R12, RDX),
  ADD_DSZ64_DRR(TMP9, TMP9, TMP15), NOP_SEQWORD },
\end{lstlisting}

The three accumulators cannot overflow.
Each unscaled limb is below $2^{51}$ and each scaled limb is below $19\cdot2^{51}$,
so every product is below $19\cdot2^{102}$ and its high half is below $2^{43}$.
Five high halves and at most five materialised carries therefore leave \code{R8}
below $2^{46}$.
Two representations of the running sum are consequently never required, and the
combined value $s_k=\code{R8}\cdot2^{64}+\code{TMP0}$ is well defined at the end of
each position.

\paragraph{Limb extraction without a mask immediate.}
At the end of a position we must split $s_k$ into a 51-bit output limb and a carry.
Immediate operands in this encoding are limited to 16 bits, so the constant
$2^{51}-1$ cannot appear in a $\mu$op and the mask cannot be applied with a single
\code{AND}.
We use two immediate shifts instead.
Shifting \code{TMP0} left by 13 and back right by 13 clears the upper 13 bits and
yields the output limb.
The same 13 bits form the low part of the carry, so we compute
$\code{TMP8}=\code{TMP0}\gg51$ and $\code{TMP1}=\code{R8}\ll13$ and combine them
with a single \code{OR} at the top of the following position, which is the
\code{OR} visible in the first triad of \Cref{lst:c25519-mulblock}.
This is exact: \code{R8} is below $2^{46}$, so $\code{R8}\ll13$ is below $2^{59}$
and no bit of the carry is lost from the 64-bit word.

These five extraction $\mu$ops do not occupy triads of their own.
They share the last three triads of the position with the final additions that
drain the high and carry accumulators, as \Cref{lst:c25519-extract} shows.
The result is that the boundary between two output positions costs no triad, in the
same way that the step boundaries of the Keccak round body merge in
\Cref{sec:cyclewalk}.

\begin{lstlisting}[caption={The closing three triads of an output position. The
extraction shifts are interleaved with the additions that drain the high and carry
accumulators, and the output limb is written in slot~1 of the last triad. Slot~2 of
that triad reads the value that slot~0 has just written to \code{R8}, which relies
on the sequential slot semantics of a triad.},label={lst:c25519-extract}]
{ ADD_DSZ64_DRR(TMP0, TMP0, RDX), SETCC_CONDB_DR(TMP15, TMP0),
  SHR_DSZ64_DRI(TMP8, TMP0, 51), NOP_SEQWORD },
{ ADD_DSZ64_DRR(R8, R8, RCX), ADD_DSZ64_DRR(TMP9, TMP9, TMP15),
  SHL_DSZ64_DRI(TMP2, TMP0, 13), NOP_SEQWORD },
{ ADD_DSZ64_DRR(R8, R8, TMP9), SHR_DSZ64_DRI(R13, TMP2, 13),
  SHL_DSZ64_DRI(TMP1, R8, 13), NOP_SEQWORD },
\end{lstlisting}

The three closing triads of the patch fold the carry leaving position four.
We multiply it by 19, add the product to limb zero, and propagate once into limb
one.
The multiply and the addition occupy adjacent slots of one triad, and the addition
reads the product that the preceding slot has just written.
Every output limb is then below $2^{51}$, except that limb one may equal $2^{51}$,
which leaves the 13 bits of headroom that the native additions and biased
subtractions of \Cref{sec:c25519-strategy} require.
This is also the bound discipline of the \texttt{amd64-51} assembly, and
\Cref{sec:c25519-verification} describes the test that holds us to it.

\subsubsection{A dedicated squaring schedule}
\label{sec:c25519-square}

Four of the nine patched operations in a ladder step are squarings, and the
inversion chain is 254 squarings and 11 multiplications.
We therefore implement squaring as a separate kernel rather than invoking
multiplication with equal operands.
Exploiting the symmetry $a_ia_j=a_ja_i$ reduces the 25 partial products to 15
distinct ones, and the resulting patch occupies 42 triads and 123 $\mu$ops at 2.93
slots per triad, against 66 triads for multiplication.
\Cref{tab:c25519-sqproducts} lists the products of each output position.

\begin{table}[htbp]
\centering
\caption{Products of the squaring patch. Each output position sums three products
instead of five, and every off-diagonal term is formed from an operand that the
wrapper has already doubled or scaled by 19. The patch itself contains no doubling
of a product and no separate reduction pass.}
\label{tab:c25519-sqproducts}
\footnotesize
\setlength{\tabcolsep}{6pt}
\begin{tabular}{@{}cl@{}}
\toprule
$k$ & products summed into $h_k$ \\
\midrule
0 & $a_0\!\cdot\!a_0$, \quad $19a_4\!\cdot\!2a_1$, \quad $19a_3\!\cdot\!2a_2$ \\
1 & $2a_0\!\cdot\!a_1$, \quad $a_3\!\cdot\!19a_3$, \quad $19a_4\!\cdot\!2a_2$ \\
2 & $2a_0\!\cdot\!a_2$, \quad $a_1\!\cdot\!a_1$, \quad $19a_4\!\cdot\!2a_3$ \\
3 & $2a_0\!\cdot\!a_3$, \quad $2a_1\!\cdot\!a_2$, \quad $a_4\!\cdot\!19a_4$ \\
4 & $2a_0\!\cdot\!a_4$, \quad $2a_1\!\cdot\!a_3$, \quad $a_2\!\cdot\!a_2$ \\
\bottomrule
\end{tabular}
\end{table}

We form the doubled and scaled operands in the wrapper rather than in the patch.
Four \code{LEA} instructions and two \code{IMUL} instructions produce
$2a_0,\ldots,2a_3$, $19a_3$ and $19a_4$ in architectural registers before the
patch is entered.
Computing them inside the patch would require six $\mu$ops and therefore two
further triads, and the combined patch budget of \Cref{sec:c25519-hooks} is the
binding constraint on this design, whereas the six native instructions are not.

The squaring patch reuses the accumulator organisation of
\Cref{sec:c25519-mul} unchanged, including the three accumulators, the
\code{SETCC} carry materialisation, the shift pair that extracts each limb, and
the final fold through 19.
Its register discipline is more aggressive, because 15 products need eight distinct
operand forms of five limbs and there are not enough registers to hold them all.
\code{RDI} holds $a_0$ on entry and receives $h_0$ at the end of position zero, at
which point $2a_0$ in \code{R15} is the only surviving form of that limb.
\code{R13} arrives holding $2a_1$, is overwritten with $a_1$ once the term
$19a_4\cdot2a_1$ has been formed, and is regenerated as $2a_1$ by a single
addition when position three requires it.
Regenerating a doubled operand from its immortal copy costs one $\mu$op, whereas
reserving a register for it would cost one of the five that receive output limbs.
Each of these rewrites occurs in a slot of the same triad that last reads the old
value, which relies on the write-after-read semantics of a triad and is the same
packing property that \Cref{sec:cyclewalk} uses for the $\pi$ cycle walk.

\subsubsection{Two co-resident patches and hook selection}
\label{sec:c25519-hooks}

Together the two kernels occupy $66+42=108$ triads and 319 $\mu$ops, and they
therefore fit simultaneously in the 128-triad patch RAM.
We place multiplication at \uaddr{7c00} and squaring immediately after it.
Each triad occupies four address units, so the squaring entry point is
$\uaddr{7c00}+66\cdot4=\uaddr{7d08}$ and the image ends at \uaddr{7db0}.
We compute the second entry point from the length of the first rather than writing
it as a constant, so that a change to the multiplication schedule cannot silently
overlap the squaring kernel.
Both addresses are even, as the match mechanism requires.

We bind multiplication to a repurposed \code{VMWRITE} path at \uaddr{0cd8} and
squaring to a repurposed \code{VMREAD} path at \uaddr{0618}.
Two properties make these instructions suitable triggers.
Both have a register-to-register form, so the triggering instruction itself
neither reads nor writes memory and requires no operand set-up that would disturb
the register interface.
Both are also unused by the surrounding program, so no other code path can reach a
patched entry point.
For squaring we emit the two opcode bytes and the ModRM byte directly, so that the
instruction presented to the decoder is exactly the form the match register
targets.

Using two hooks rather than one is what makes the dedicated squaring schedule
usable.
A single hook would require rewriting a match register between a multiplication and
a squaring, and the ladder alternates between them within every step.
With both hooks installed once at start-up, each of the 2561 invocations of a
scalar multiplication reaches its kernel directly.
The saturated control of \Cref{sec:c25519-challenge} illustrates the alternative:
its multiplier alone occupies 75 triads, leaving no room for a co-resident squaring
kernel, so it must compute every square as a multiplication with equal operands.

\subsubsection{Register chaining across field operations}
\label{sec:c25519-integration}

A conventional C interface reloads the input limbs before every field operation and
stores the output limbs immediately afterwards.
We retain the stores, because later consumers read the value from memory, but we
avoid the reload whenever the next operation consumes the value just produced.
We place one complete ladder step in a single inline-assembly block, use \code{RBP}
as the base of a fixed-offset structure holding all fifteen field elements of the
step, and address every operand as a constant displacement from it.
Additions and subtractions are then written so that their five results are left in
\code{RDI}, \code{RSI}, \code{R12}, \code{R11} and \code{R14}, which is exactly the
first-operand register set of both patches.

The interface of \Cref{tab:c25519-regif} determines where this is possible.
An addition or subtraction is free to choose its output registers, so it can always
be chained into a patched operation.
The output set of squaring shares only \code{RDI} with the first-operand set, so two
squarings can be chained at a cost of four register renames.
The output set of multiplication shares no register with the first-operand set, so
two multiplications cannot be chained and the second reloads from memory.
\Cref{tab:c25519-chains} lists the seven producer--consumer pairs in a ladder step
where the reload is removed.

\begin{table}[htbp]
\centering
\caption{Producer--consumer pairs in one ladder step for which the consumer reads
its first operand from registers rather than memory. Each pair removes five loads.
The producer is a native addition or subtraction in every case, since the output
register set of a patched multiplication does not match its own first-operand set.}
\label{tab:c25519-chains}
\footnotesize
\setlength{\tabcolsep}{6pt}
\begin{tabular}{@{}lll@{}}
\toprule
Producer & Consumer & Result \\
\midrule
$A=x_2+z_2$          & squaring          & $AA$      \\
$B=x_2-z_2$          & squaring          & $BB$      \\
$D=x_3-z_3$          & multiplication by $A$ & $DA$  \\
$C=x_3+z_3$          & multiplication by $B$ & $CB$  \\
$DA+CB$              & squaring          & new $x_3$ \\
$DA-CB$              & squaring          & $(DA-CB)^2$ \\
$AA+121665E$         & multiplication by $E$ & new $z_2$ \\
\bottomrule
\end{tabular}
\end{table}

Obtaining the fourth pair requires one change to the order of the step.
The natural order computes $C$, then $D$, then $DA$, then $CB$, which chains only
$D$ into $DA$.
We compute $D$ and $DA$ first and then $C$ and $CB$, which chains both.
The two orders are equivalent because $C$ and $D$ depend only on $x_3$ and $z_3$ and
feed different multiplications.

The inversion contains far longer chains of a single operation.
We evaluate $z^{p-2}$ with a fixed addition chain whose runs of consecutive
squarings have lengths 1, 2, 1, 5, 10, 20, 10, 50, 100, 50 and 5, separated by 11
multiplications.
Within a run we rename the four output registers that differ from the input set and
invoke the patch again, so every intermediate of the run stays in registers.
Only the first squaring of a run loads a field element and only the last stores one,
which removes 486 of the 508 field-element transfers that the same chain would perform
through a per-operation interface.
We emit the whole inversion as one assembly block, because the register state must
survive from the beginning of a run to its end and a function boundary would not
permit that.

\subsubsection{Patch layout and verification}
\label{sec:c25519-verification}

We validate the installed patches rather than a model of them.
The wrappers declare every register that a patch clobbers, and each wrapper keeps
its state or output pointer in a register that the patch leaves intact, so that a
missing clobber declaration would be observed as a wrong result rather than as a
silent miscompilation.

We test the complete X25519 implementation rather than the two patched operations
in isolation.
The tests comprise the two standalone vectors of RFC~7748, the single-iteration
vector, and the 1000-iteration chain, and they exercise scalar clamping, all 255
ladder steps, the inversion, the final reduction and the encoding.
We run the identical tests for the hand-written C, fiat-crypto, CryptOpt,
\texttt{amd64-51} and \texttt{amd64-64} contenders of \Cref{sec:c25519-perf}, and we
compare the microcode output against each native backend after the chain has
completed.

The 1000-iteration chain also serves a second purpose.
It is our gate on the limb bound discipline described at the end of
\Cref{sec:c25519-mul}.
Each iteration feeds the encoded output of the previous one back as the next scalar
and point, so an output limb that exceeded its permitted range would propagate into
the following scalar multiplication and diverge from the reference within a few
iterations rather than producing a result that happens to agree on one input.
The same chain gates the substitution of the \texttt{amd64-51} assembly field
operations into our ladder, which is only sound if the two implementations agree on
how much headroom a returned limb may consume.

Neither patch contains a conditional operation or an indexed memory access.
Every operand is a register named in the patch image, every shift amount is an
immediate fixed at installation time, and the schedule of each output position is
identical regardless of operand values.
The surrounding ladder performs its conditional swap arithmetically and iterates a
fixed 255 times.
No data-dependent branch or data-dependent address therefore arises, and as in
\Cref{sec:keccak} we assume that the underlying $\mu$ops have data-independent
latency.
 in place of the existing
% subsection, or paste inline.
%
% Every triad, uop and slot-density figure below is taken from the installed
% patch arrays in simple/curve25519/full_curve25519_inline2.c.
%
% NOTE: the original text cited \Cref{sec:ucode-layer}, which is not defined
% anywhere in gold489.tex. This version states the intra-triad forwarding
% property locally and cites \Cref{sec:cyclewalk}, which establishes the same
% slot semantics for Keccak. When \section{Background and Programming Model}
% (sec:background) is written, that claim probably belongs there, cited from
% both case studies.

\subsection{Implementation}

\subsubsection{Register interface and operand lifetimes}
\label{sec:c25519-interface}

Neither patch performs a memory access.
All operands, accumulators, carries, and results live in registers for the whole
invocation, and the surrounding x86 wrapper is responsible for placing the inputs
and retrieving the outputs.
This division follows from the resource that binds each side.
Patch RAM is scarce and native instruction slots are not, so we move every
operation that does not require the microcode interface into the wrapper.
\Cref{tab:c25519-regif} gives the resulting register interface.

The wrapper also supplies two pre-zeroed accumulator seeds in \code{RAX} and
\code{R8}, which removes two initialisation $\mu$ops from each patch.
For multiplication the wrapper issues ten loads and two register clears.
For squaring it issues five loads, four \code{LEA} instructions that form the
doubled operands, two \code{IMUL} instructions that form the operands scaled by
19, and the same two clears.

\begin{table}[htbp]
\centering
\caption{Register interface of the two field-arithmetic patches.
The first operand of multiplication occupies five registers that no
$\mu$op overwrites, so all 25 partial products read them in place.
The second operand is copied to \code{TMP10}--\code{TMP14} at patch entry, which
frees its five architectural registers to receive output limbs as successive
positions complete.
The squaring patch receives the doubled and scaled operands from the wrapper and
returns its result in a different register set, which \Cref{sec:c25519-integration}
uses to explain where chaining is possible.}
\label{tab:c25519-regif}
\footnotesize
\setlength{\tabcolsep}{4pt}
\begin{tabularx}{\linewidth}{@{}llX@{}}
\toprule
Role & Registers & Lifetime \\
\midrule
\multicolumn{3}{@{}l}{\emph{multiplication}} \\
$a_0,\ldots,a_4$      & \code{RDI RSI R12 R11 R14} & read-only for the entire patch \\
$b_0,\ldots,b_4$      & \code{R15 R13 R9 R10 RBX}  & copied to \code{TMP10}--\code{TMP14}, then reused as outputs \\
$19b_1,\ldots,19b_4$  & \code{R13 R9 R10 RBX}      & formed in place at entry, consumed in decreasing index \\
$h_0,\ldots,h_4$      & \code{R15 R13 R9 R10 RAX}  & written as each output position completes \\
\midrule
\multicolumn{3}{@{}l}{\emph{squaring}} \\
$a_0,\ldots,a_4$      & \code{RDI RSI R12 R11 R14} & \code{RDI} is overwritten by $h_0$ at the end of position zero \\
$2a_0,\ldots,2a_3$    & \code{R15 R13 R9 R10}      & supplied by the wrapper; \code{R13} is regenerated inside the patch \\
$19a_4$, $19a_3$      & \code{RBX RDX}             & supplied by the wrapper \\
$h_0,\ldots,h_4$      & \code{RDI R9 R10 RBX RAX}  & written as each output position completes \\
\midrule
\multicolumn{3}{@{}l}{\emph{shared internal state}} \\
low accumulator       & \code{TMP0}   & holds the incoming carry, then the sum of the low product halves \\
high accumulator      & \code{R8}     & sum of the high product halves \\
carry accumulator     & \code{TMP9}   & count of carries out of the low accumulator \\
materialised carry    & \code{TMP15}  & destination of \code{SETCC\_CONDB} \\
product high half     & \code{RCX}    & second output of every multiply \\
product staging       & \code{RDX}    & the operand each multiply is allowed to destroy \\
limb extraction       & \code{TMP1 TMP2 TMP8} & shift results combined into the next carry \\
\bottomrule
\end{tabularx}
\end{table}

The unusual feature of this interface is the treatment of the second operand,
and it follows from the shape of the multiply $\mu$op.
\code{MUL\_DSZ64\_DRR} has two destinations.
It preserves its first source, writes the low half of the product over its
\emph{second} source, and writes the high half to the named destination.
Every product therefore destroys one of the operands that formed it.
A schoolbook multiplication reads each limb of both operands five times, so no
limb can be left in the register that a product will consume.

We resolve this by separating an immortal copy of each operand from a mortal
staging register.
The first operand is immortal by construction, since it always occupies the
preserved source position.
The second operand is copied to \code{TMP10}--\code{TMP14} in the first triad of
the patch, and thereafter each product is formed by copying the required limb into
\code{RDX} and letting the multiply overwrite it.
The three-triad prologue performs this copy and forms $19b_1,\ldots,19b_4$ in
place in the architectural registers, using the same destructive low-half write
that the schoolbook products use.
Both forms of the second operand are consequently live at once, the unscaled form
in temporary registers and the scaled form in architectural ones.

\subsubsection{Register-resident multiplication}
\label{sec:c25519-mul}

The multiplication patch occupies 66 triads and contains 196 $\mu$ops, an average
of 2.97 of the three available slots.
It is straight-line code with no microsequencer control flow.
This distinguishes it from the Keccak kernel of \Cref{sec:keccak}, where a loop
was unavoidable because a full permutation cannot be represented in 128 triads.
A field operation contains no repetition worth keeping resident, and its complete
body fits, so we spend no triad on loop control and no $\mu$op on a loop variable.

The patch computes the five output positions in sequence.
Three triads form the prologue, twelve triads compute each output position, and
three triads perform the final reduction.
For position $k$ we accumulate the five products required by the recurrence of
\Cref{sec:c25519-strategy}, extract the low 51 bits as limb $k$, and pass the
remaining bits to position $k+1$.

\paragraph{The product staircase.}
The five twelve-triad blocks are the same schedule with different operands.
This regularity is not incidental, and \Cref{tab:c25519-staircase} shows where it
comes from.
The first operand advances through $a_0,\ldots,a_4$ in the same order in every
block, so its five registers are read in a fixed sequence.
Only the second operand sequence moves, and it moves by exactly one position per
block: it begins at $b_k$, descends through the unscaled limbs to $b_0$, and then
continues into the scaled limbs from $19b_4$ downwards.
Position four requires no scaled term, because $i+j=9$ has no solution with
$i,j\leq4$, and position three requires only the single term $19a_4b_4$.

\begin{table}[htbp]
\centering
\caption{Second-operand sequence for each output position of the multiplication
patch. The first operand is always read as $a_0,a_1,a_2,a_3,a_4$ in that order,
so each row lists only the register staged into \code{RDX} for the corresponding
product. Unscaled limbs are read from the temporary copies and scaled limbs from
the architectural registers. The rightmost column names the register that
receives the completed output limb, which in every case is the architectural
register whose scaled multiple was last required by the preceding row.}
\label{tab:c25519-staircase}
\footnotesize
\setlength{\tabcolsep}{5pt}
\begin{tabular}{@{}clllllc@{}}
\toprule
$k$ & $\times a_0$ & $\times a_1$ & $\times a_2$ & $\times a_3$ & $\times a_4$ & $h_k$ \\
\midrule
0 & $b_0$ & $19b_4$ & $19b_3$ & $19b_2$ & $19b_1$ & \code{R15} \\
1 & $b_1$ & $b_0$    & $19b_4$ & $19b_3$ & $19b_2$ & \code{R13} \\
2 & $b_2$ & $b_1$    & $b_0$    & $19b_4$ & $19b_3$ & \code{R9}  \\
3 & $b_3$ & $b_2$    & $b_1$    & $b_0$    & $19b_4$ & \code{R10} \\
4 & $b_4$ & $b_3$    & $b_2$    & $b_1$    & $b_0$    & \code{RAX} \\
\bottomrule
\end{tabular}
\end{table}

The staircase determines where the output limbs are placed.
The scaled multiple $19b_j$ is required for the last time in the block for
position $4-j$, and the register holding it is dead from that point onward.
We write the output limb of that block into precisely that register.
Limb zero therefore lands in \code{R15}, which held $b_0$ and whose value is
preserved in \code{TMP10}, and limbs one through three land in \code{R13},
\code{R9} and \code{R10} as $19b_1$, $19b_2$ and $19b_3$ retire in turn.
Only limb four requires a register that never held part of an operand, and it
uses \code{RAX}, which the wrapper supplies as a zeroed seed and which is free once
the low accumulator has been initialised.
The patch thus completes a full multiplication without ever holding more than one
copy of any operand limb and without a single spill.

\paragraph{Three accumulators and materialised carries.}
Each output position sums five 128-bit products.
We represent the running sum in three registers rather than two.
\code{TMP0} accumulates the low halves, \code{R8} accumulates the high halves, and
\code{TMP9} counts the carries that leave \code{TMP0}.

The third accumulator is what the microcode carry interface buys us.
On Goldmont, each arithmetic $\mu$op records its carry with its destination
register rather than in a single architectural flag, and
\code{SETCC\_CONDB\_DR} converts that recorded carry into a zero or one value.
We place the \code{SETCC} in the slot immediately after the addition that produces
the carry, so the two occupy one triad and the materialisation costs no additional
triad.
The materialised value is then added to \code{TMP9} while the high-half chain
proceeds in a different slot.
A conventional \code{ADC} chain would instead serialise the two halves, because
\code{ADC} reads the architectural carry flag and only one carry can be pending in
that flag at a time.
Goldmont implements neither ADX nor a second carry flag, so native code for this
representation has no equivalent of the third accumulator.
Keeping the carry count separate costs one triad per output position and removes a
serialisation point from the dependency path of all five.

Two properties of the microcode interface constrain how this chain is written.
\code{SETCC\_CONDB\_DR} writes a correct value only when its destination is a
temporary register, so all three accumulators and the materialised carry are
temporary registers and the chain never touches architectural state.
The recorded carry also persists only until the next arithmetic operation writes
the same destination, which fixes the position of the \code{SETCC} relative to its
addition and is the reason the two are packed together rather than scheduled
apart.

\begin{lstlisting}[caption={The opening five triads of an output position, taken
from position one. Slot~2 of the first triad forms the incoming carry, slot~0 of
each even triad stages the next second-operand limb into \code{RDX}, and each
addition into \code{TMP0} is immediately followed by the \code{SETCC} that
materialises its carry. The high accumulator and the carry accumulator are
initialised in the third triad, after which both chains advance in parallel with
the multiplies.},label={lst:c25519-mulblock}]
{ ZEROEXT_DSZ64_DR(RDX, TMP11), MUL_DSZ64_DRR(RCX, RDI, RDX),
  OR_DSZ64_DRR(TMP0, TMP8, TMP1), NOP_SEQWORD },
{ ADD_DSZ64_DRR(TMP0, TMP0, RDX), SETCC_CONDB_DR(TMP15, TMP0),
  ZEROEXT_DSZ64_DR(RDX, TMP10), NOP_SEQWORD },
{ ZEROEXT_DSZ64_DR(R8, RCX), MUL_DSZ64_DRR(RCX, RSI, RDX),
  ZEROEXT_DSZ64_DR(TMP9, TMP15), NOP_SEQWORD },
{ ADD_DSZ64_DRR(TMP0, TMP0, RDX), SETCC_CONDB_DR(TMP15, TMP0),
  ZEROEXT_DSZ64_DR(RDX, RBX), NOP_SEQWORD },
{ ADD_DSZ64_DRR(R8, R8, RCX), MUL_DSZ64_DRR(RCX, R12, RDX),
  ADD_DSZ64_DRR(TMP9, TMP9, TMP15), NOP_SEQWORD },
\end{lstlisting}

The three accumulators cannot overflow.
Each unscaled limb is below $2^{51}$ and each scaled limb is below $19\cdot2^{51}$,
so every product is below $19\cdot2^{102}$ and its high half is below $2^{43}$.
Five high halves and at most five materialised carries therefore leave \code{R8}
below $2^{46}$.
Two representations of the running sum are consequently never required, and the
combined value $s_k=\code{R8}\cdot2^{64}+\code{TMP0}$ is well defined at the end of
each position.

\paragraph{Limb extraction without a mask immediate.}
At the end of a position we must split $s_k$ into a 51-bit output limb and a carry.
Immediate operands in this encoding are limited to 16 bits, so the constant
$2^{51}-1$ cannot appear in a $\mu$op and the mask cannot be applied with a single
\code{AND}.
We use two immediate shifts instead.
Shifting \code{TMP0} left by 13 and back right by 13 clears the upper 13 bits and
yields the output limb.
The same 13 bits form the low part of the carry, so we compute
$\code{TMP8}=\code{TMP0}\gg51$ and $\code{TMP1}=\code{R8}\ll13$ and combine them
with a single \code{OR} at the top of the following position, which is the
\code{OR} visible in the first triad of \Cref{lst:c25519-mulblock}.
This is exact: \code{R8} is below $2^{46}$, so $\code{R8}\ll13$ is below $2^{59}$
and no bit of the carry is lost from the 64-bit word.

These five extraction $\mu$ops do not occupy triads of their own.
They share the last three triads of the position with the final additions that
drain the high and carry accumulators, as \Cref{lst:c25519-extract} shows.
The result is that the boundary between two output positions costs no triad, in the
same way that the step boundaries of the Keccak round body merge in
\Cref{sec:cyclewalk}.

\begin{lstlisting}[caption={The closing three triads of an output position. The
extraction shifts are interleaved with the additions that drain the high and carry
accumulators, and the output limb is written in slot~1 of the last triad. Slot~2 of
that triad reads the value that slot~0 has just written to \code{R8}, which relies
on the sequential slot semantics of a triad.},label={lst:c25519-extract}]
{ ADD_DSZ64_DRR(TMP0, TMP0, RDX), SETCC_CONDB_DR(TMP15, TMP0),
  SHR_DSZ64_DRI(TMP8, TMP0, 51), NOP_SEQWORD },
{ ADD_DSZ64_DRR(R8, R8, RCX), ADD_DSZ64_DRR(TMP9, TMP9, TMP15),
  SHL_DSZ64_DRI(TMP2, TMP0, 13), NOP_SEQWORD },
{ ADD_DSZ64_DRR(R8, R8, TMP9), SHR_DSZ64_DRI(R13, TMP2, 13),
  SHL_DSZ64_DRI(TMP1, R8, 13), NOP_SEQWORD },
\end{lstlisting}

The three closing triads of the patch fold the carry leaving position four.
We multiply it by 19, add the product to limb zero, and propagate once into limb
one.
The multiply and the addition occupy adjacent slots of one triad, and the addition
reads the product that the preceding slot has just written.
Every output limb is then below $2^{51}$, except that limb one may equal $2^{51}$,
which leaves the 13 bits of headroom that the native additions and biased
subtractions of \Cref{sec:c25519-strategy} require.
This is also the bound discipline of the \texttt{amd64-51} assembly, and
\Cref{sec:c25519-verification} describes the test that holds us to it.

\subsubsection{A dedicated squaring schedule}
\label{sec:c25519-square}

Four of the nine patched operations in a ladder step are squarings, and the
inversion chain is 254 squarings and 11 multiplications.
We therefore implement squaring as a separate kernel rather than invoking
multiplication with equal operands.
Exploiting the symmetry $a_ia_j=a_ja_i$ reduces the 25 partial products to 15
distinct ones, and the resulting patch occupies 42 triads and 123 $\mu$ops at 2.93
slots per triad, against 66 triads for multiplication.
\Cref{tab:c25519-sqproducts} lists the products of each output position.

\begin{table}[htbp]
\centering
\caption{Products of the squaring patch. Each output position sums three products
instead of five, and every off-diagonal term is formed from an operand that the
wrapper has already doubled or scaled by 19. The patch itself contains no doubling
of a product and no separate reduction pass.}
\label{tab:c25519-sqproducts}
\footnotesize
\setlength{\tabcolsep}{6pt}
\begin{tabular}{@{}cl@{}}
\toprule
$k$ & products summed into $h_k$ \\
\midrule
0 & $a_0\!\cdot\!a_0$, \quad $19a_4\!\cdot\!2a_1$, \quad $19a_3\!\cdot\!2a_2$ \\
1 & $2a_0\!\cdot\!a_1$, \quad $a_3\!\cdot\!19a_3$, \quad $19a_4\!\cdot\!2a_2$ \\
2 & $2a_0\!\cdot\!a_2$, \quad $a_1\!\cdot\!a_1$, \quad $19a_4\!\cdot\!2a_3$ \\
3 & $2a_0\!\cdot\!a_3$, \quad $2a_1\!\cdot\!a_2$, \quad $a_4\!\cdot\!19a_4$ \\
4 & $2a_0\!\cdot\!a_4$, \quad $2a_1\!\cdot\!a_3$, \quad $a_2\!\cdot\!a_2$ \\
\bottomrule
\end{tabular}
\end{table}

We form the doubled and scaled operands in the wrapper rather than in the patch.
Four \code{LEA} instructions and two \code{IMUL} instructions produce
$2a_0,\ldots,2a_3$, $19a_3$ and $19a_4$ in architectural registers before the
patch is entered.
Computing them inside the patch would require six $\mu$ops and therefore two
further triads, and the combined patch budget of \Cref{sec:c25519-hooks} is the
binding constraint on this design, whereas the six native instructions are not.

The squaring patch reuses the accumulator organisation of
\Cref{sec:c25519-mul} unchanged, including the three accumulators, the
\code{SETCC} carry materialisation, the shift pair that extracts each limb, and
the final fold through 19.
Its register discipline is more aggressive, because 15 products need eight distinct
operand forms of five limbs and there are not enough registers to hold them all.
\code{RDI} holds $a_0$ on entry and receives $h_0$ at the end of position zero, at
which point $2a_0$ in \code{R15} is the only surviving form of that limb.
\code{R13} arrives holding $2a_1$, is overwritten with $a_1$ once the term
$19a_4\cdot2a_1$ has been formed, and is regenerated as $2a_1$ by a single
addition when position three requires it.
Regenerating a doubled operand from its immortal copy costs one $\mu$op, whereas
reserving a register for it would cost one of the five that receive output limbs.
Each of these rewrites occurs in a slot of the same triad that last reads the old
value, which relies on the write-after-read semantics of a triad and is the same
packing property that \Cref{sec:cyclewalk} uses for the $\pi$ cycle walk.

\subsubsection{Two co-resident patches and hook selection}
\label{sec:c25519-hooks}

Together the two kernels occupy $66+42=108$ triads and 319 $\mu$ops, and they
therefore fit simultaneously in the 128-triad patch RAM.
We place multiplication at \uaddr{7c00} and squaring immediately after it.
Each triad occupies four address units, so the squaring entry point is
$\uaddr{7c00}+66\cdot4=\uaddr{7d08}$ and the image ends at \uaddr{7db0}.
We compute the second entry point from the length of the first rather than writing
it as a constant, so that a change to the multiplication schedule cannot silently
overlap the squaring kernel.
Both addresses are even, as the match mechanism requires.

We bind multiplication to a repurposed \code{VMWRITE} path at \uaddr{0cd8} and
squaring to a repurposed \code{VMREAD} path at \uaddr{0618}.
Two properties make these instructions suitable triggers.
Both have a register-to-register form, so the triggering instruction itself
neither reads nor writes memory and requires no operand set-up that would disturb
the register interface.
Both are also unused by the surrounding program, so no other code path can reach a
patched entry point.
For squaring we emit the two opcode bytes and the ModRM byte directly, so that the
instruction presented to the decoder is exactly the form the match register
targets.

Using two hooks rather than one is what makes the dedicated squaring schedule
usable.
A single hook would require rewriting a match register between a multiplication and
a squaring, and the ladder alternates between them within every step.
With both hooks installed once at start-up, each of the 2561 invocations of a
scalar multiplication reaches its kernel directly.
The saturated control of \Cref{sec:c25519-challenge} illustrates the alternative:
its multiplier alone occupies 75 triads, leaving no room for a co-resident squaring
kernel, so it must compute every square as a multiplication with equal operands.

\subsubsection{Register chaining across field operations}
\label{sec:c25519-integration}

A conventional C interface reloads the input limbs before every field operation and
stores the output limbs immediately afterwards.
We retain the stores, because later consumers read the value from memory, but we
avoid the reload whenever the next operation consumes the value just produced.
We place one complete ladder step in a single inline-assembly block, use \code{RBP}
as the base of a fixed-offset structure holding all fifteen field elements of the
step, and address every operand as a constant displacement from it.
Additions and subtractions are then written so that their five results are left in
\code{RDI}, \code{RSI}, \code{R12}, \code{R11} and \code{R14}, which is exactly the
first-operand register set of both patches.

The interface of \Cref{tab:c25519-regif} determines where this is possible.
An addition or subtraction is free to choose its output registers, so it can always
be chained into a patched operation.
The output set of squaring shares only \code{RDI} with the first-operand set, so two
squarings can be chained at a cost of four register renames.
The output set of multiplication shares no register with the first-operand set, so
two multiplications cannot be chained and the second reloads from memory.
\Cref{tab:c25519-chains} lists the seven producer--consumer pairs in a ladder step
where the reload is removed.

\begin{table}[htbp]
\centering
\caption{Producer--consumer pairs in one ladder step for which the consumer reads
its first operand from registers rather than memory. Each pair removes five loads.
The producer is a native addition or subtraction in every case, since the output
register set of a patched multiplication does not match its own first-operand set.}
\label{tab:c25519-chains}
\footnotesize
\setlength{\tabcolsep}{6pt}
\begin{tabular}{@{}lll@{}}
\toprule
Producer & Consumer & Result \\
\midrule
$A=x_2+z_2$          & squaring          & $AA$      \\
$B=x_2-z_2$          & squaring          & $BB$      \\
$D=x_3-z_3$          & multiplication by $A$ & $DA$  \\
$C=x_3+z_3$          & multiplication by $B$ & $CB$  \\
$DA+CB$              & squaring          & new $x_3$ \\
$DA-CB$              & squaring          & $(DA-CB)^2$ \\
$AA+121665E$         & multiplication by $E$ & new $z_2$ \\
\bottomrule
\end{tabular}
\end{table}

Obtaining the fourth pair requires one change to the order of the step.
The natural order computes $C$, then $D$, then $DA$, then $CB$, which chains only
$D$ into $DA$.
We compute $D$ and $DA$ first and then $C$ and $CB$, which chains both.
The two orders are equivalent because $C$ and $D$ depend only on $x_3$ and $z_3$ and
feed different multiplications.

The inversion contains far longer chains of a single operation.
We evaluate $z^{p-2}$ with a fixed addition chain whose runs of consecutive
squarings have lengths 1, 2, 1, 5, 10, 20, 10, 50, 100, 50 and 5, separated by 11
multiplications.
Within a run we rename the four output registers that differ from the input set and
invoke the patch again, so every intermediate of the run stays in registers.
Only the first squaring of a run loads a field element and only the last stores one,
which removes 486 of the 508 field-element transfers that the same chain would perform
through a per-operation interface.
We emit the whole inversion as one assembly block, because the register state must
survive from the beginning of a run to its end and a function boundary would not
permit that.

\subsubsection{Patch layout and verification}
\label{sec:c25519-verification}

We validate the installed patches rather than a model of them.
The wrappers declare every register that a patch clobbers, and each wrapper keeps
its state or output pointer in a register that the patch leaves intact, so that a
missing clobber declaration would be observed as a wrong result rather than as a
silent miscompilation.

We test the complete X25519 implementation rather than the two patched operations
in isolation.
The tests comprise the two standalone vectors of RFC~7748, the single-iteration
vector, and the 1000-iteration chain, and they exercise scalar clamping, all 255
ladder steps, the inversion, the final reduction and the encoding.
We run the identical tests for the hand-written C, fiat-crypto, CryptOpt,
\texttt{amd64-51} and \texttt{amd64-64} contenders of \Cref{sec:c25519-perf}, and we
compare the microcode output against each native backend after the chain has
completed.

The 1000-iteration chain also serves a second purpose.
It is our gate on the limb bound discipline described at the end of
\Cref{sec:c25519-mul}.
Each iteration feeds the encoded output of the previous one back as the next scalar
and point, so an output limb that exceeded its permitted range would propagate into
the following scalar multiplication and diverge from the reference within a few
iterations rather than producing a result that happens to agree on one input.
The same chain gates the substitution of the \texttt{amd64-51} assembly field
operations into our ladder, which is only sound if the two implementations agree on
how much headroom a returned limb may consume.

Neither patch contains a conditional operation or an indexed memory access.
Every operand is a register named in the patch image, every shift amount is an
immediate fixed at installation time, and the schedule of each output position is
identical regardless of operand values.
The surrounding ladder performs its conditional swap arithmetically and iterates a
fixed 255 times.
No data-dependent branch or data-dependent address therefore arises, and as in
\Cref{sec:keccak} we assume that the underlying $\mu$ops have data-independent
latency.
 in place of the existing
% subsection, or paste inline.
%
% Every triad, uop and slot-density figure below is taken from the installed
% patch arrays in simple/curve25519/full_curve25519_inline2.c.
%
% NOTE: the original text cited \Cref{sec:ucode-layer}, which is not defined
% anywhere in gold489.tex. This version states the intra-triad forwarding
% property locally and cites \Cref{sec:cyclewalk}, which establishes the same
% slot semantics for Keccak. When \section{Background and Programming Model}
% (sec:background) is written, that claim probably belongs there, cited from
% both case studies.

\subsection{Implementation}

\subsubsection{Register interface and operand lifetimes}
\label{sec:c25519-interface}

Neither patch performs a memory access.
All operands, accumulators, carries, and results live in registers for the whole
invocation, and the surrounding x86 wrapper is responsible for placing the inputs
and retrieving the outputs.
This division follows from the resource that binds each side.
Patch RAM is scarce and native instruction slots are not, so we move every
operation that does not require the microcode interface into the wrapper.
\Cref{tab:c25519-regif} gives the resulting register interface.

The wrapper also supplies two pre-zeroed accumulator seeds in \code{RAX} and
\code{R8}, which removes two initialisation $\mu$ops from each patch.
For multiplication the wrapper issues ten loads and two register clears.
For squaring it issues five loads, four \code{LEA} instructions that form the
doubled operands, two \code{IMUL} instructions that form the operands scaled by
19, and the same two clears.

\begin{table}[htbp]
\centering
\caption{Register interface of the two field-arithmetic patches.
The first operand of multiplication occupies five registers that no
$\mu$op overwrites, so all 25 partial products read them in place.
The second operand is copied to \code{TMP10}--\code{TMP14} at patch entry, which
frees its five architectural registers to receive output limbs as successive
positions complete.
The squaring patch receives the doubled and scaled operands from the wrapper and
returns its result in a different register set, which \Cref{sec:c25519-integration}
uses to explain where chaining is possible.}
\label{tab:c25519-regif}
\footnotesize
\setlength{\tabcolsep}{4pt}
\begin{tabularx}{\linewidth}{@{}llX@{}}
\toprule
Role & Registers & Lifetime \\
\midrule
\multicolumn{3}{@{}l}{\emph{multiplication}} \\
$a_0,\ldots,a_4$      & \code{RDI RSI R12 R11 R14} & read-only for the entire patch \\
$b_0,\ldots,b_4$      & \code{R15 R13 R9 R10 RBX}  & copied to \code{TMP10}--\code{TMP14}, then reused as outputs \\
$19b_1,\ldots,19b_4$  & \code{R13 R9 R10 RBX}      & formed in place at entry, consumed in decreasing index \\
$h_0,\ldots,h_4$      & \code{R15 R13 R9 R10 RAX}  & written as each output position completes \\
\midrule
\multicolumn{3}{@{}l}{\emph{squaring}} \\
$a_0,\ldots,a_4$      & \code{RDI RSI R12 R11 R14} & \code{RDI} is overwritten by $h_0$ at the end of position zero \\
$2a_0,\ldots,2a_3$    & \code{R15 R13 R9 R10}      & supplied by the wrapper; \code{R13} is regenerated inside the patch \\
$19a_4$, $19a_3$      & \code{RBX RDX}             & supplied by the wrapper \\
$h_0,\ldots,h_4$      & \code{RDI R9 R10 RBX RAX}  & written as each output position completes \\
\midrule
\multicolumn{3}{@{}l}{\emph{shared internal state}} \\
low accumulator       & \code{TMP0}   & holds the incoming carry, then the sum of the low product halves \\
high accumulator      & \code{R8}     & sum of the high product halves \\
carry accumulator     & \code{TMP9}   & count of carries out of the low accumulator \\
materialised carry    & \code{TMP15}  & destination of \code{SETCC\_CONDB} \\
product high half     & \code{RCX}    & second output of every multiply \\
product staging       & \code{RDX}    & the operand each multiply is allowed to destroy \\
limb extraction       & \code{TMP1 TMP2 TMP8} & shift results combined into the next carry \\
\bottomrule
\end{tabularx}
\end{table}

The unusual feature of this interface is the treatment of the second operand,
and it follows from the shape of the multiply $\mu$op.
\code{MUL\_DSZ64\_DRR} has two destinations.
It preserves its first source, writes the low half of the product over its
\emph{second} source, and writes the high half to the named destination.
Every product therefore destroys one of the operands that formed it.
A schoolbook multiplication reads each limb of both operands five times, so no
limb can be left in the register that a product will consume.

We resolve this by separating an immortal copy of each operand from a mortal
staging register.
The first operand is immortal by construction, since it always occupies the
preserved source position.
The second operand is copied to \code{TMP10}--\code{TMP14} in the first triad of
the patch, and thereafter each product is formed by copying the required limb into
\code{RDX} and letting the multiply overwrite it.
The three-triad prologue performs this copy and forms $19b_1,\ldots,19b_4$ in
place in the architectural registers, using the same destructive low-half write
that the schoolbook products use.
Both forms of the second operand are consequently live at once, the unscaled form
in temporary registers and the scaled form in architectural ones.

\subsubsection{Register-resident multiplication}
\label{sec:c25519-mul}

The multiplication patch occupies 66 triads and contains 196 $\mu$ops, an average
of 2.97 of the three available slots.
It is straight-line code with no microsequencer control flow.
This distinguishes it from the Keccak kernel of \Cref{sec:keccak}, where a loop
was unavoidable because a full permutation cannot be represented in 128 triads.
A field operation contains no repetition worth keeping resident, and its complete
body fits, so we spend no triad on loop control and no $\mu$op on a loop variable.

The patch computes the five output positions in sequence.
Three triads form the prologue, twelve triads compute each output position, and
three triads perform the final reduction.
For position $k$ we accumulate the five products required by the recurrence of
\Cref{sec:c25519-strategy}, extract the low 51 bits as limb $k$, and pass the
remaining bits to position $k+1$.

\paragraph{The product staircase.}
The five twelve-triad blocks are the same schedule with different operands.
This regularity is not incidental, and \Cref{tab:c25519-staircase} shows where it
comes from.
The first operand advances through $a_0,\ldots,a_4$ in the same order in every
block, so its five registers are read in a fixed sequence.
Only the second operand sequence moves, and it moves by exactly one position per
block: it begins at $b_k$, descends through the unscaled limbs to $b_0$, and then
continues into the scaled limbs from $19b_4$ downwards.
Position four requires no scaled term, because $i+j=9$ has no solution with
$i,j\leq4$, and position three requires only the single term $19a_4b_4$.

\begin{table}[htbp]
\centering
\caption{Second-operand sequence for each output position of the multiplication
patch. The first operand is always read as $a_0,a_1,a_2,a_3,a_4$ in that order,
so each row lists only the register staged into \code{RDX} for the corresponding
product. Unscaled limbs are read from the temporary copies and scaled limbs from
the architectural registers. The rightmost column names the register that
receives the completed output limb, which in every case is the architectural
register whose scaled multiple was last required by the preceding row.}
\label{tab:c25519-staircase}
\footnotesize
\setlength{\tabcolsep}{5pt}
\begin{tabular}{@{}clllllc@{}}
\toprule
$k$ & $\times a_0$ & $\times a_1$ & $\times a_2$ & $\times a_3$ & $\times a_4$ & $h_k$ \\
\midrule
0 & $b_0$ & $19b_4$ & $19b_3$ & $19b_2$ & $19b_1$ & \code{R15} \\
1 & $b_1$ & $b_0$    & $19b_4$ & $19b_3$ & $19b_2$ & \code{R13} \\
2 & $b_2$ & $b_1$    & $b_0$    & $19b_4$ & $19b_3$ & \code{R9}  \\
3 & $b_3$ & $b_2$    & $b_1$    & $b_0$    & $19b_4$ & \code{R10} \\
4 & $b_4$ & $b_3$    & $b_2$    & $b_1$    & $b_0$    & \code{RAX} \\
\bottomrule
\end{tabular}
\end{table}

The staircase determines where the output limbs are placed.
The scaled multiple $19b_j$ is required for the last time in the block for
position $4-j$, and the register holding it is dead from that point onward.
We write the output limb of that block into precisely that register.
Limb zero therefore lands in \code{R15}, which held $b_0$ and whose value is
preserved in \code{TMP10}, and limbs one through three land in \code{R13},
\code{R9} and \code{R10} as $19b_1$, $19b_2$ and $19b_3$ retire in turn.
Only limb four requires a register that never held part of an operand, and it
uses \code{RAX}, which the wrapper supplies as a zeroed seed and which is free once
the low accumulator has been initialised.
The patch thus completes a full multiplication without ever holding more than one
copy of any operand limb and without a single spill.

\paragraph{Three accumulators and materialised carries.}
Each output position sums five 128-bit products.
We represent the running sum in three registers rather than two.
\code{TMP0} accumulates the low halves, \code{R8} accumulates the high halves, and
\code{TMP9} counts the carries that leave \code{TMP0}.

The third accumulator is what the microcode carry interface buys us.
On Goldmont, each arithmetic $\mu$op records its carry with its destination
register rather than in a single architectural flag, and
\code{SETCC\_CONDB\_DR} converts that recorded carry into a zero or one value.
We place the \code{SETCC} in the slot immediately after the addition that produces
the carry, so the two occupy one triad and the materialisation costs no additional
triad.
The materialised value is then added to \code{TMP9} while the high-half chain
proceeds in a different slot.
A conventional \code{ADC} chain would instead serialise the two halves, because
\code{ADC} reads the architectural carry flag and only one carry can be pending in
that flag at a time.
Goldmont implements neither ADX nor a second carry flag, so native code for this
representation has no equivalent of the third accumulator.
Keeping the carry count separate costs one triad per output position and removes a
serialisation point from the dependency path of all five.

Two properties of the microcode interface constrain how this chain is written.
\code{SETCC\_CONDB\_DR} writes a correct value only when its destination is a
temporary register, so all three accumulators and the materialised carry are
temporary registers and the chain never touches architectural state.
The recorded carry also persists only until the next arithmetic operation writes
the same destination, which fixes the position of the \code{SETCC} relative to its
addition and is the reason the two are packed together rather than scheduled
apart.

\begin{lstlisting}[caption={The opening five triads of an output position, taken
from position one. Slot~2 of the first triad forms the incoming carry, slot~0 of
each even triad stages the next second-operand limb into \code{RDX}, and each
addition into \code{TMP0} is immediately followed by the \code{SETCC} that
materialises its carry. The high accumulator and the carry accumulator are
initialised in the third triad, after which both chains advance in parallel with
the multiplies.},label={lst:c25519-mulblock}]
{ ZEROEXT_DSZ64_DR(RDX, TMP11), MUL_DSZ64_DRR(RCX, RDI, RDX),
  OR_DSZ64_DRR(TMP0, TMP8, TMP1), NOP_SEQWORD },
{ ADD_DSZ64_DRR(TMP0, TMP0, RDX), SETCC_CONDB_DR(TMP15, TMP0),
  ZEROEXT_DSZ64_DR(RDX, TMP10), NOP_SEQWORD },
{ ZEROEXT_DSZ64_DR(R8, RCX), MUL_DSZ64_DRR(RCX, RSI, RDX),
  ZEROEXT_DSZ64_DR(TMP9, TMP15), NOP_SEQWORD },
{ ADD_DSZ64_DRR(TMP0, TMP0, RDX), SETCC_CONDB_DR(TMP15, TMP0),
  ZEROEXT_DSZ64_DR(RDX, RBX), NOP_SEQWORD },
{ ADD_DSZ64_DRR(R8, R8, RCX), MUL_DSZ64_DRR(RCX, R12, RDX),
  ADD_DSZ64_DRR(TMP9, TMP9, TMP15), NOP_SEQWORD },
\end{lstlisting}

The three accumulators cannot overflow.
Each unscaled limb is below $2^{51}$ and each scaled limb is below $19\cdot2^{51}$,
so every product is below $19\cdot2^{102}$ and its high half is below $2^{43}$.
Five high halves and at most five materialised carries therefore leave \code{R8}
below $2^{46}$.
Two representations of the running sum are consequently never required, and the
combined value $s_k=\code{R8}\cdot2^{64}+\code{TMP0}$ is well defined at the end of
each position.

\paragraph{Limb extraction without a mask immediate.}
At the end of a position we must split $s_k$ into a 51-bit output limb and a carry.
Immediate operands in this encoding are limited to 16 bits, so the constant
$2^{51}-1$ cannot appear in a $\mu$op and the mask cannot be applied with a single
\code{AND}.
We use two immediate shifts instead.
Shifting \code{TMP0} left by 13 and back right by 13 clears the upper 13 bits and
yields the output limb.
The same 13 bits form the low part of the carry, so we compute
$\code{TMP8}=\code{TMP0}\gg51$ and $\code{TMP1}=\code{R8}\ll13$ and combine them
with a single \code{OR} at the top of the following position, which is the
\code{OR} visible in the first triad of \Cref{lst:c25519-mulblock}.
This is exact: \code{R8} is below $2^{46}$, so $\code{R8}\ll13$ is below $2^{59}$
and no bit of the carry is lost from the 64-bit word.

These five extraction $\mu$ops do not occupy triads of their own.
They share the last three triads of the position with the final additions that
drain the high and carry accumulators, as \Cref{lst:c25519-extract} shows.
The result is that the boundary between two output positions costs no triad, in the
same way that the step boundaries of the Keccak round body merge in
\Cref{sec:cyclewalk}.

\begin{lstlisting}[caption={The closing three triads of an output position. The
extraction shifts are interleaved with the additions that drain the high and carry
accumulators, and the output limb is written in slot~1 of the last triad. Slot~2 of
that triad reads the value that slot~0 has just written to \code{R8}, which relies
on the sequential slot semantics of a triad.},label={lst:c25519-extract}]
{ ADD_DSZ64_DRR(TMP0, TMP0, RDX), SETCC_CONDB_DR(TMP15, TMP0),
  SHR_DSZ64_DRI(TMP8, TMP0, 51), NOP_SEQWORD },
{ ADD_DSZ64_DRR(R8, R8, RCX), ADD_DSZ64_DRR(TMP9, TMP9, TMP15),
  SHL_DSZ64_DRI(TMP2, TMP0, 13), NOP_SEQWORD },
{ ADD_DSZ64_DRR(R8, R8, TMP9), SHR_DSZ64_DRI(R13, TMP2, 13),
  SHL_DSZ64_DRI(TMP1, R8, 13), NOP_SEQWORD },
\end{lstlisting}

The three closing triads of the patch fold the carry leaving position four.
We multiply it by 19, add the product to limb zero, and propagate once into limb
one.
The multiply and the addition occupy adjacent slots of one triad, and the addition
reads the product that the preceding slot has just written.
Every output limb is then below $2^{51}$, except that limb one may equal $2^{51}$,
which leaves the 13 bits of headroom that the native additions and biased
subtractions of \Cref{sec:c25519-strategy} require.
This is also the bound discipline of the \texttt{amd64-51} assembly, and
\Cref{sec:c25519-verification} describes the test that holds us to it.

\subsubsection{A dedicated squaring schedule}
\label{sec:c25519-square}

Four of the nine patched operations in a ladder step are squarings, and the
inversion chain is 254 squarings and 11 multiplications.
We therefore implement squaring as a separate kernel rather than invoking
multiplication with equal operands.
Exploiting the symmetry $a_ia_j=a_ja_i$ reduces the 25 partial products to 15
distinct ones, and the resulting patch occupies 42 triads and 123 $\mu$ops at 2.93
slots per triad, against 66 triads for multiplication.
\Cref{tab:c25519-sqproducts} lists the products of each output position.

\begin{table}[htbp]
\centering
\caption{Products of the squaring patch. Each output position sums three products
instead of five, and every off-diagonal term is formed from an operand that the
wrapper has already doubled or scaled by 19. The patch itself contains no doubling
of a product and no separate reduction pass.}
\label{tab:c25519-sqproducts}
\footnotesize
\setlength{\tabcolsep}{6pt}
\begin{tabular}{@{}cl@{}}
\toprule
$k$ & products summed into $h_k$ \\
\midrule
0 & $a_0\!\cdot\!a_0$, \quad $19a_4\!\cdot\!2a_1$, \quad $19a_3\!\cdot\!2a_2$ \\
1 & $2a_0\!\cdot\!a_1$, \quad $a_3\!\cdot\!19a_3$, \quad $19a_4\!\cdot\!2a_2$ \\
2 & $2a_0\!\cdot\!a_2$, \quad $a_1\!\cdot\!a_1$, \quad $19a_4\!\cdot\!2a_3$ \\
3 & $2a_0\!\cdot\!a_3$, \quad $2a_1\!\cdot\!a_2$, \quad $a_4\!\cdot\!19a_4$ \\
4 & $2a_0\!\cdot\!a_4$, \quad $2a_1\!\cdot\!a_3$, \quad $a_2\!\cdot\!a_2$ \\
\bottomrule
\end{tabular}
\end{table}

We form the doubled and scaled operands in the wrapper rather than in the patch.
Four \code{LEA} instructions and two \code{IMUL} instructions produce
$2a_0,\ldots,2a_3$, $19a_3$ and $19a_4$ in architectural registers before the
patch is entered.
Computing them inside the patch would require six $\mu$ops and therefore two
further triads, and the combined patch budget of \Cref{sec:c25519-hooks} is the
binding constraint on this design, whereas the six native instructions are not.

The squaring patch reuses the accumulator organisation of
\Cref{sec:c25519-mul} unchanged, including the three accumulators, the
\code{SETCC} carry materialisation, the shift pair that extracts each limb, and
the final fold through 19.
Its register discipline is more aggressive, because 15 products need eight distinct
operand forms of five limbs and there are not enough registers to hold them all.
\code{RDI} holds $a_0$ on entry and receives $h_0$ at the end of position zero, at
which point $2a_0$ in \code{R15} is the only surviving form of that limb.
\code{R13} arrives holding $2a_1$, is overwritten with $a_1$ once the term
$19a_4\cdot2a_1$ has been formed, and is regenerated as $2a_1$ by a single
addition when position three requires it.
Regenerating a doubled operand from its immortal copy costs one $\mu$op, whereas
reserving a register for it would cost one of the five that receive output limbs.
Each of these rewrites occurs in a slot of the same triad that last reads the old
value, which relies on the write-after-read semantics of a triad and is the same
packing property that \Cref{sec:cyclewalk} uses for the $\pi$ cycle walk.

\subsubsection{Two co-resident patches and hook selection}
\label{sec:c25519-hooks}

Together the two kernels occupy $66+42=108$ triads and 319 $\mu$ops, and they
therefore fit simultaneously in the 128-triad patch RAM.
We place multiplication at \uaddr{7c00} and squaring immediately after it.
Each triad occupies four address units, so the squaring entry point is
$\uaddr{7c00}+66\cdot4=\uaddr{7d08}$ and the image ends at \uaddr{7db0}.
We compute the second entry point from the length of the first rather than writing
it as a constant, so that a change to the multiplication schedule cannot silently
overlap the squaring kernel.
Both addresses are even, as the match mechanism requires.

We bind multiplication to a repurposed \code{VMWRITE} path at \uaddr{0cd8} and
squaring to a repurposed \code{VMREAD} path at \uaddr{0618}.
Two properties make these instructions suitable triggers.
Both have a register-to-register form, so the triggering instruction itself
neither reads nor writes memory and requires no operand set-up that would disturb
the register interface.
Both are also unused by the surrounding program, so no other code path can reach a
patched entry point.
For squaring we emit the two opcode bytes and the ModRM byte directly, so that the
instruction presented to the decoder is exactly the form the match register
targets.

Using two hooks rather than one is what makes the dedicated squaring schedule
usable.
A single hook would require rewriting a match register between a multiplication and
a squaring, and the ladder alternates between them within every step.
With both hooks installed once at start-up, each of the 2561 invocations of a
scalar multiplication reaches its kernel directly.
The saturated control of \Cref{sec:c25519-challenge} illustrates the alternative:
its multiplier alone occupies 75 triads, leaving no room for a co-resident squaring
kernel, so it must compute every square as a multiplication with equal operands.

\subsubsection{Register chaining across field operations}
\label{sec:c25519-integration}

A conventional C interface reloads the input limbs before every field operation and
stores the output limbs immediately afterwards.
We retain the stores, because later consumers read the value from memory, but we
avoid the reload whenever the next operation consumes the value just produced.
We place one complete ladder step in a single inline-assembly block, use \code{RBP}
as the base of a fixed-offset structure holding all fifteen field elements of the
step, and address every operand as a constant displacement from it.
Additions and subtractions are then written so that their five results are left in
\code{RDI}, \code{RSI}, \code{R12}, \code{R11} and \code{R14}, which is exactly the
first-operand register set of both patches.

The interface of \Cref{tab:c25519-regif} determines where this is possible.
An addition or subtraction is free to choose its output registers, so it can always
be chained into a patched operation.
The output set of squaring shares only \code{RDI} with the first-operand set, so two
squarings can be chained at a cost of four register renames.
The output set of multiplication shares no register with the first-operand set, so
two multiplications cannot be chained and the second reloads from memory.
\Cref{tab:c25519-chains} lists the seven producer--consumer pairs in a ladder step
where the reload is removed.

\begin{table}[htbp]
\centering
\caption{Producer--consumer pairs in one ladder step for which the consumer reads
its first operand from registers rather than memory. Each pair removes five loads.
The producer is a native addition or subtraction in every case, since the output
register set of a patched multiplication does not match its own first-operand set.}
\label{tab:c25519-chains}
\footnotesize
\setlength{\tabcolsep}{6pt}
\begin{tabular}{@{}lll@{}}
\toprule
Producer & Consumer & Result \\
\midrule
$A=x_2+z_2$          & squaring          & $AA$      \\
$B=x_2-z_2$          & squaring          & $BB$      \\
$D=x_3-z_3$          & multiplication by $A$ & $DA$  \\
$C=x_3+z_3$          & multiplication by $B$ & $CB$  \\
$DA+CB$              & squaring          & new $x_3$ \\
$DA-CB$              & squaring          & $(DA-CB)^2$ \\
$AA+121665E$         & multiplication by $E$ & new $z_2$ \\
\bottomrule
\end{tabular}
\end{table}

Obtaining the fourth pair requires one change to the order of the step.
The natural order computes $C$, then $D$, then $DA$, then $CB$, which chains only
$D$ into $DA$.
We compute $D$ and $DA$ first and then $C$ and $CB$, which chains both.
The two orders are equivalent because $C$ and $D$ depend only on $x_3$ and $z_3$ and
feed different multiplications.

The inversion contains far longer chains of a single operation.
We evaluate $z^{p-2}$ with a fixed addition chain whose runs of consecutive
squarings have lengths 1, 2, 1, 5, 10, 20, 10, 50, 100, 50 and 5, separated by 11
multiplications.
Within a run we rename the four output registers that differ from the input set and
invoke the patch again, so every intermediate of the run stays in registers.
Only the first squaring of a run loads a field element and only the last stores one,
which removes 486 of the 508 field-element transfers that the same chain would perform
through a per-operation interface.
We emit the whole inversion as one assembly block, because the register state must
survive from the beginning of a run to its end and a function boundary would not
permit that.

\subsubsection{Patch layout and verification}
\label{sec:c25519-verification}

We validate the installed patches rather than a model of them.
The wrappers declare every register that a patch clobbers, and each wrapper keeps
its state or output pointer in a register that the patch leaves intact, so that a
missing clobber declaration would be observed as a wrong result rather than as a
silent miscompilation.

We test the complete X25519 implementation rather than the two patched operations
in isolation.
The tests comprise the two standalone vectors of RFC~7748, the single-iteration
vector, and the 1000-iteration chain, and they exercise scalar clamping, all 255
ladder steps, the inversion, the final reduction and the encoding.
We run the identical tests for the hand-written C, fiat-crypto, CryptOpt,
\texttt{amd64-51} and \texttt{amd64-64} contenders of \Cref{sec:c25519-perf}, and we
compare the microcode output against each native backend after the chain has
completed.

The 1000-iteration chain also serves a second purpose.
It is our gate on the limb bound discipline described at the end of
\Cref{sec:c25519-mul}.
Each iteration feeds the encoded output of the previous one back as the next scalar
and point, so an output limb that exceeded its permitted range would propagate into
the following scalar multiplication and diverge from the reference within a few
iterations rather than producing a result that happens to agree on one input.
The same chain gates the substitution of the \texttt{amd64-51} assembly field
operations into our ladder, which is only sound if the two implementations agree on
how much headroom a returned limb may consume.

Neither patch contains a conditional operation or an indexed memory access.
Every operand is a register named in the patch image, every shift amount is an
immediate fixed at installation time, and the schedule of each output position is
identical regardless of operand values.
The surrounding ladder performs its conditional swap arithmetically and iterates a
fixed 255 times.
No data-dependent branch or data-dependent address therefore arises, and as in
\Cref{sec:keccak} we assume that the underlying $\mu$ops have data-independent
latency.
 in place of the existing
% subsection, or paste inline.
%
% Every triad, uop and slot-density figure below is taken from the installed
% patch arrays in simple/curve25519/full_curve25519_inline2.c.
%
% NOTE: the original text cited \Cref{sec:ucode-layer}, which is not defined
% anywhere in gold489.tex. This version states the intra-triad forwarding
% property locally and cites \Cref{sec:cyclewalk}, which establishes the same
% slot semantics for Keccak. When \section{Background and Programming Model}
% (sec:background) is written, that claim probably belongs there, cited from
% both case studies.

\subsection{Implementation}

\subsubsection{Register interface and operand lifetimes}
\label{sec:c25519-interface}

Neither patch performs a memory access.
All operands, accumulators, carries, and results live in registers for the whole
invocation, and the surrounding x86 wrapper is responsible for placing the inputs
and retrieving the outputs.
This division follows from the resource that binds each side.
Patch RAM is scarce and native instruction slots are not, so we move every
operation that does not require the microcode interface into the wrapper.
\Cref{tab:c25519-regif} gives the resulting register interface.

The wrapper also supplies two pre-zeroed accumulator seeds in \code{RAX} and
\code{R8}, which removes two initialisation $\mu$ops from each patch.
For multiplication the wrapper issues ten loads and two register clears.
For squaring it issues five loads, four \code{LEA} instructions that form the
doubled operands, two \code{IMUL} instructions that form the operands scaled by
19, and the same two clears.

\begin{table}[htbp]
\centering
\caption{Register interface of the two field-arithmetic patches.
The first operand of multiplication occupies five registers that no
$\mu$op overwrites, so all 25 partial products read them in place.
The second operand is copied to \code{TMP10}--\code{TMP14} at patch entry, which
frees its five architectural registers to receive output limbs as successive
positions complete.
The squaring patch receives the doubled and scaled operands from the wrapper and
returns its result in a different register set, which \Cref{sec:c25519-integration}
uses to explain where chaining is possible.}
\label{tab:c25519-regif}
\footnotesize
\setlength{\tabcolsep}{4pt}
\begin{tabularx}{\linewidth}{@{}llX@{}}
\toprule
Role & Registers & Lifetime \\
\midrule
\multicolumn{3}{@{}l}{\emph{multiplication}} \\
$a_0,\ldots,a_4$      & \code{RDI RSI R12 R11 R14} & read-only for the entire patch \\
$b_0,\ldots,b_4$      & \code{R15 R13 R9 R10 RBX}  & copied to \code{TMP10}--\code{TMP14}, then reused as outputs \\
$19b_1,\ldots,19b_4$  & \code{R13 R9 R10 RBX}      & formed in place at entry, consumed in decreasing index \\
$h_0,\ldots,h_4$      & \code{R15 R13 R9 R10 RAX}  & written as each output position completes \\
\midrule
\multicolumn{3}{@{}l}{\emph{squaring}} \\
$a_0,\ldots,a_4$      & \code{RDI RSI R12 R11 R14} & \code{RDI} is overwritten by $h_0$ at the end of position zero \\
$2a_0,\ldots,2a_3$    & \code{R15 R13 R9 R10}      & supplied by the wrapper; \code{R13} is regenerated inside the patch \\
$19a_4$, $19a_3$      & \code{RBX RDX}             & supplied by the wrapper \\
$h_0,\ldots,h_4$      & \code{RDI R9 R10 RBX RAX}  & written as each output position completes \\
\midrule
\multicolumn{3}{@{}l}{\emph{shared internal state}} \\
low accumulator       & \code{TMP0}   & holds the incoming carry, then the sum of the low product halves \\
high accumulator      & \code{R8}     & sum of the high product halves \\
carry accumulator     & \code{TMP9}   & count of carries out of the low accumulator \\
materialised carry    & \code{TMP15}  & destination of \code{SETCC\_CONDB} \\
product high half     & \code{RCX}    & second output of every multiply \\
product staging       & \code{RDX}    & the operand each multiply is allowed to destroy \\
limb extraction       & \code{TMP1 TMP2 TMP8} & shift results combined into the next carry \\
\bottomrule
\end{tabularx}
\end{table}

The unusual feature of this interface is the treatment of the second operand,
and it follows from the shape of the multiply $\mu$op.
\code{MUL\_DSZ64\_DRR} has two destinations.
It preserves its first source, writes the low half of the product over its
\emph{second} source, and writes the high half to the named destination.
Every product therefore destroys one of the operands that formed it.
A schoolbook multiplication reads each limb of both operands five times, so no
limb can be left in the register that a product will consume.

We resolve this by separating an immortal copy of each operand from a mortal
staging register.
The first operand is immortal by construction, since it always occupies the
preserved source position.
The second operand is copied to \code{TMP10}--\code{TMP14} in the first triad of
the patch, and thereafter each product is formed by copying the required limb into
\code{RDX} and letting the multiply overwrite it.
The three-triad prologue performs this copy and forms $19b_1,\ldots,19b_4$ in
place in the architectural registers, using the same destructive low-half write
that the schoolbook products use.
Both forms of the second operand are consequently live at once, the unscaled form
in temporary registers and the scaled form in architectural ones.

\subsubsection{Register-resident multiplication}
\label{sec:c25519-mul}

The multiplication patch occupies 66 triads and contains 196 $\mu$ops, an average
of 2.97 of the three available slots.
It is straight-line code with no microsequencer control flow.
This distinguishes it from the Keccak kernel of \Cref{sec:keccak}, where a loop
was unavoidable because a full permutation cannot be represented in 128 triads.
A field operation contains no repetition worth keeping resident, and its complete
body fits, so we spend no triad on loop control and no $\mu$op on a loop variable.

The patch computes the five output positions in sequence.
Three triads form the prologue, twelve triads compute each output position, and
three triads perform the final reduction.
For position $k$ we accumulate the five products required by the recurrence of
\Cref{sec:c25519-strategy}, extract the low 51 bits as limb $k$, and pass the
remaining bits to position $k+1$.

\paragraph{The product staircase.}
The five twelve-triad blocks are the same schedule with different operands.
This regularity is not incidental, and \Cref{tab:c25519-staircase} shows where it
comes from.
The first operand advances through $a_0,\ldots,a_4$ in the same order in every
block, so its five registers are read in a fixed sequence.
Only the second operand sequence moves, and it moves by exactly one position per
block: it begins at $b_k$, descends through the unscaled limbs to $b_0$, and then
continues into the scaled limbs from $19b_4$ downwards.
Position four requires no scaled term, because $i+j=9$ has no solution with
$i,j\leq4$, and position three requires only the single term $19a_4b_4$.

\begin{table}[htbp]
\centering
\caption{Second-operand sequence for each output position of the multiplication
patch. The first operand is always read as $a_0,a_1,a_2,a_3,a_4$ in that order,
so each row lists only the register staged into \code{RDX} for the corresponding
product. Unscaled limbs are read from the temporary copies and scaled limbs from
the architectural registers. The rightmost column names the register that
receives the completed output limb, which in every case is the architectural
register whose scaled multiple was last required by the preceding row.}
\label{tab:c25519-staircase}
\footnotesize
\setlength{\tabcolsep}{5pt}
\begin{tabular}{@{}clllllc@{}}
\toprule
$k$ & $\times a_0$ & $\times a_1$ & $\times a_2$ & $\times a_3$ & $\times a_4$ & $h_k$ \\
\midrule
0 & $b_0$ & $19b_4$ & $19b_3$ & $19b_2$ & $19b_1$ & \code{R15} \\
1 & $b_1$ & $b_0$    & $19b_4$ & $19b_3$ & $19b_2$ & \code{R13} \\
2 & $b_2$ & $b_1$    & $b_0$    & $19b_4$ & $19b_3$ & \code{R9}  \\
3 & $b_3$ & $b_2$    & $b_1$    & $b_0$    & $19b_4$ & \code{R10} \\
4 & $b_4$ & $b_3$    & $b_2$    & $b_1$    & $b_0$    & \code{RAX} \\
\bottomrule
\end{tabular}
\end{table}

The staircase determines where the output limbs are placed.
The scaled multiple $19b_j$ is required for the last time in the block for
position $4-j$, and the register holding it is dead from that point onward.
We write the output limb of that block into precisely that register.
Limb zero therefore lands in \code{R15}, which held $b_0$ and whose value is
preserved in \code{TMP10}, and limbs one through three land in \code{R13},
\code{R9} and \code{R10} as $19b_1$, $19b_2$ and $19b_3$ retire in turn.
Only limb four requires a register that never held part of an operand, and it
uses \code{RAX}, which the wrapper supplies as a zeroed seed and which is free once
the low accumulator has been initialised.
The patch thus completes a full multiplication without ever holding more than one
copy of any operand limb and without a single spill.

\paragraph{Three accumulators and materialised carries.}
Each output position sums five 128-bit products.
We represent the running sum in three registers rather than two.
\code{TMP0} accumulates the low halves, \code{R8} accumulates the high halves, and
\code{TMP9} counts the carries that leave \code{TMP0}.

The third accumulator is what the microcode carry interface buys us.
On Goldmont, each arithmetic $\mu$op records its carry with its destination
register rather than in a single architectural flag, and
\code{SETCC\_CONDB\_DR} converts that recorded carry into a zero or one value.
We place the \code{SETCC} in the slot immediately after the addition that produces
the carry, so the two occupy one triad and the materialisation costs no additional
triad.
The materialised value is then added to \code{TMP9} while the high-half chain
proceeds in a different slot.
A conventional \code{ADC} chain would instead serialise the two halves, because
\code{ADC} reads the architectural carry flag and only one carry can be pending in
that flag at a time.
Goldmont implements neither ADX nor a second carry flag, so native code for this
representation has no equivalent of the third accumulator.
Keeping the carry count separate costs one triad per output position and removes a
serialisation point from the dependency path of all five.

Two properties of the microcode interface constrain how this chain is written.
\code{SETCC\_CONDB\_DR} writes a correct value only when its destination is a
temporary register, so all three accumulators and the materialised carry are
temporary registers and the chain never touches architectural state.
The recorded carry also persists only until the next arithmetic operation writes
the same destination, which fixes the position of the \code{SETCC} relative to its
addition and is the reason the two are packed together rather than scheduled
apart.

\begin{lstlisting}[caption={The opening five triads of an output position, taken
from position one. Slot~2 of the first triad forms the incoming carry, slot~0 of
each even triad stages the next second-operand limb into \code{RDX}, and each
addition into \code{TMP0} is immediately followed by the \code{SETCC} that
materialises its carry. The high accumulator and the carry accumulator are
initialised in the third triad, after which both chains advance in parallel with
the multiplies.},label={lst:c25519-mulblock}]
{ ZEROEXT_DSZ64_DR(RDX, TMP11), MUL_DSZ64_DRR(RCX, RDI, RDX),
  OR_DSZ64_DRR(TMP0, TMP8, TMP1), NOP_SEQWORD },
{ ADD_DSZ64_DRR(TMP0, TMP0, RDX), SETCC_CONDB_DR(TMP15, TMP0),
  ZEROEXT_DSZ64_DR(RDX, TMP10), NOP_SEQWORD },
{ ZEROEXT_DSZ64_DR(R8, RCX), MUL_DSZ64_DRR(RCX, RSI, RDX),
  ZEROEXT_DSZ64_DR(TMP9, TMP15), NOP_SEQWORD },
{ ADD_DSZ64_DRR(TMP0, TMP0, RDX), SETCC_CONDB_DR(TMP15, TMP0),
  ZEROEXT_DSZ64_DR(RDX, RBX), NOP_SEQWORD },
{ ADD_DSZ64_DRR(R8, R8, RCX), MUL_DSZ64_DRR(RCX, R12, RDX),
  ADD_DSZ64_DRR(TMP9, TMP9, TMP15), NOP_SEQWORD },
\end{lstlisting}

The three accumulators cannot overflow.
Each unscaled limb is below $2^{51}$ and each scaled limb is below $19\cdot2^{51}$,
so every product is below $19\cdot2^{102}$ and its high half is below $2^{43}$.
Five high halves and at most five materialised carries therefore leave \code{R8}
below $2^{46}$.
Two representations of the running sum are consequently never required, and the
combined value $s_k=\code{R8}\cdot2^{64}+\code{TMP0}$ is well defined at the end of
each position.

\paragraph{Limb extraction without a mask immediate.}
At the end of a position we must split $s_k$ into a 51-bit output limb and a carry.
Immediate operands in this encoding are limited to 16 bits, so the constant
$2^{51}-1$ cannot appear in a $\mu$op and the mask cannot be applied with a single
\code{AND}.
We use two immediate shifts instead.
Shifting \code{TMP0} left by 13 and back right by 13 clears the upper 13 bits and
yields the output limb.
The same 13 bits form the low part of the carry, so we compute
$\code{TMP8}=\code{TMP0}\gg51$ and $\code{TMP1}=\code{R8}\ll13$ and combine them
with a single \code{OR} at the top of the following position, which is the
\code{OR} visible in the first triad of \Cref{lst:c25519-mulblock}.
This is exact: \code{R8} is below $2^{46}$, so $\code{R8}\ll13$ is below $2^{59}$
and no bit of the carry is lost from the 64-bit word.

These five extraction $\mu$ops do not occupy triads of their own.
They share the last three triads of the position with the final additions that
drain the high and carry accumulators, as \Cref{lst:c25519-extract} shows.
The result is that the boundary between two output positions costs no triad, in the
same way that the step boundaries of the Keccak round body merge in
\Cref{sec:cyclewalk}.

\begin{lstlisting}[caption={The closing three triads of an output position. The
extraction shifts are interleaved with the additions that drain the high and carry
accumulators, and the output limb is written in slot~1 of the last triad. Slot~2 of
that triad reads the value that slot~0 has just written to \code{R8}, which relies
on the sequential slot semantics of a triad.},label={lst:c25519-extract}]
{ ADD_DSZ64_DRR(TMP0, TMP0, RDX), SETCC_CONDB_DR(TMP15, TMP0),
  SHR_DSZ64_DRI(TMP8, TMP0, 51), NOP_SEQWORD },
{ ADD_DSZ64_DRR(R8, R8, RCX), ADD_DSZ64_DRR(TMP9, TMP9, TMP15),
  SHL_DSZ64_DRI(TMP2, TMP0, 13), NOP_SEQWORD },
{ ADD_DSZ64_DRR(R8, R8, TMP9), SHR_DSZ64_DRI(R13, TMP2, 13),
  SHL_DSZ64_DRI(TMP1, R8, 13), NOP_SEQWORD },
\end{lstlisting}

The three closing triads of the patch fold the carry leaving position four.
We multiply it by 19, add the product to limb zero, and propagate once into limb
one.
The multiply and the addition occupy adjacent slots of one triad, and the addition
reads the product that the preceding slot has just written.
Every output limb is then below $2^{51}$, except that limb one may equal $2^{51}$,
which leaves the 13 bits of headroom that the native additions and biased
subtractions of \Cref{sec:c25519-strategy} require.
This is also the bound discipline of the \texttt{amd64-51} assembly, and
\Cref{sec:c25519-verification} describes the test that holds us to it.

\subsubsection{A dedicated squaring schedule}
\label{sec:c25519-square}

Four of the nine patched operations in a ladder step are squarings, and the
inversion chain is 254 squarings and 11 multiplications.
We therefore implement squaring as a separate kernel rather than invoking
multiplication with equal operands.
Exploiting the symmetry $a_ia_j=a_ja_i$ reduces the 25 partial products to 15
distinct ones, and the resulting patch occupies 42 triads and 123 $\mu$ops at 2.93
slots per triad, against 66 triads for multiplication.
\Cref{tab:c25519-sqproducts} lists the products of each output position.

\begin{table}[htbp]
\centering
\caption{Products of the squaring patch. Each output position sums three products
instead of five, and every off-diagonal term is formed from an operand that the
wrapper has already doubled or scaled by 19. The patch itself contains no doubling
of a product and no separate reduction pass.}
\label{tab:c25519-sqproducts}
\footnotesize
\setlength{\tabcolsep}{6pt}
\begin{tabular}{@{}cl@{}}
\toprule
$k$ & products summed into $h_k$ \\
\midrule
0 & $a_0\!\cdot\!a_0$, \quad $19a_4\!\cdot\!2a_1$, \quad $19a_3\!\cdot\!2a_2$ \\
1 & $2a_0\!\cdot\!a_1$, \quad $a_3\!\cdot\!19a_3$, \quad $19a_4\!\cdot\!2a_2$ \\
2 & $2a_0\!\cdot\!a_2$, \quad $a_1\!\cdot\!a_1$, \quad $19a_4\!\cdot\!2a_3$ \\
3 & $2a_0\!\cdot\!a_3$, \quad $2a_1\!\cdot\!a_2$, \quad $a_4\!\cdot\!19a_4$ \\
4 & $2a_0\!\cdot\!a_4$, \quad $2a_1\!\cdot\!a_3$, \quad $a_2\!\cdot\!a_2$ \\
\bottomrule
\end{tabular}
\end{table}

We form the doubled and scaled operands in the wrapper rather than in the patch.
Four \code{LEA} instructions and two \code{IMUL} instructions produce
$2a_0,\ldots,2a_3$, $19a_3$ and $19a_4$ in architectural registers before the
patch is entered.
Computing them inside the patch would require six $\mu$ops and therefore two
further triads, and the combined patch budget of \Cref{sec:c25519-hooks} is the
binding constraint on this design, whereas the six native instructions are not.

The squaring patch reuses the accumulator organisation of
\Cref{sec:c25519-mul} unchanged, including the three accumulators, the
\code{SETCC} carry materialisation, the shift pair that extracts each limb, and
the final fold through 19.
Its register discipline is more aggressive, because 15 products need eight distinct
operand forms of five limbs and there are not enough registers to hold them all.
\code{RDI} holds $a_0$ on entry and receives $h_0$ at the end of position zero, at
which point $2a_0$ in \code{R15} is the only surviving form of that limb.
\code{R13} arrives holding $2a_1$, is overwritten with $a_1$ once the term
$19a_4\cdot2a_1$ has been formed, and is regenerated as $2a_1$ by a single
addition when position three requires it.
Regenerating a doubled operand from its immortal copy costs one $\mu$op, whereas
reserving a register for it would cost one of the five that receive output limbs.
Each of these rewrites occurs in a slot of the same triad that last reads the old
value, which relies on the write-after-read semantics of a triad and is the same
packing property that \Cref{sec:cyclewalk} uses for the $\pi$ cycle walk.

\subsubsection{Two co-resident patches and hook selection}
\label{sec:c25519-hooks}

Together the two kernels occupy $66+42=108$ triads and 319 $\mu$ops, and they
therefore fit simultaneously in the 128-triad patch RAM.
We place multiplication at \uaddr{7c00} and squaring immediately after it.
Each triad occupies four address units, so the squaring entry point is
$\uaddr{7c00}+66\cdot4=\uaddr{7d08}$ and the image ends at \uaddr{7db0}.
We compute the second entry point from the length of the first rather than writing
it as a constant, so that a change to the multiplication schedule cannot silently
overlap the squaring kernel.
Both addresses are even, as the match mechanism requires.

We bind multiplication to a repurposed \code{VMWRITE} path at \uaddr{0cd8} and
squaring to a repurposed \code{VMREAD} path at \uaddr{0618}.
Two properties make these instructions suitable triggers.
Both have a register-to-register form, so the triggering instruction itself
neither reads nor writes memory and requires no operand set-up that would disturb
the register interface.
Both are also unused by the surrounding program, so no other code path can reach a
patched entry point.
For squaring we emit the two opcode bytes and the ModRM byte directly, so that the
instruction presented to the decoder is exactly the form the match register
targets.

Using two hooks rather than one is what makes the dedicated squaring schedule
usable.
A single hook would require rewriting a match register between a multiplication and
a squaring, and the ladder alternates between them within every step.
With both hooks installed once at start-up, each of the 2561 invocations of a
scalar multiplication reaches its kernel directly.
The saturated control of \Cref{sec:c25519-challenge} illustrates the alternative:
its multiplier alone occupies 75 triads, leaving no room for a co-resident squaring
kernel, so it must compute every square as a multiplication with equal operands.

\subsubsection{Register chaining across field operations}
\label{sec:c25519-integration}

A conventional C interface reloads the input limbs before every field operation and
stores the output limbs immediately afterwards.
We retain the stores, because later consumers read the value from memory, but we
avoid the reload whenever the next operation consumes the value just produced.
We place one complete ladder step in a single inline-assembly block, use \code{RBP}
as the base of a fixed-offset structure holding all fifteen field elements of the
step, and address every operand as a constant displacement from it.
Additions and subtractions are then written so that their five results are left in
\code{RDI}, \code{RSI}, \code{R12}, \code{R11} and \code{R14}, which is exactly the
first-operand register set of both patches.

The interface of \Cref{tab:c25519-regif} determines where this is possible.
An addition or subtraction is free to choose its output registers, so it can always
be chained into a patched operation.
The output set of squaring shares only \code{RDI} with the first-operand set, so two
squarings can be chained at a cost of four register renames.
The output set of multiplication shares no register with the first-operand set, so
two multiplications cannot be chained and the second reloads from memory.
\Cref{tab:c25519-chains} lists the seven producer--consumer pairs in a ladder step
where the reload is removed.

\begin{table}[htbp]
\centering
\caption{Producer--consumer pairs in one ladder step for which the consumer reads
its first operand from registers rather than memory. Each pair removes five loads.
The producer is a native addition or subtraction in every case, since the output
register set of a patched multiplication does not match its own first-operand set.}
\label{tab:c25519-chains}
\footnotesize
\setlength{\tabcolsep}{6pt}
\begin{tabular}{@{}lll@{}}
\toprule
Producer & Consumer & Result \\
\midrule
$A=x_2+z_2$          & squaring          & $AA$      \\
$B=x_2-z_2$          & squaring          & $BB$      \\
$D=x_3-z_3$          & multiplication by $A$ & $DA$  \\
$C=x_3+z_3$          & multiplication by $B$ & $CB$  \\
$DA+CB$              & squaring          & new $x_3$ \\
$DA-CB$              & squaring          & $(DA-CB)^2$ \\
$AA+121665E$         & multiplication by $E$ & new $z_2$ \\
\bottomrule
\end{tabular}
\end{table}

Obtaining the fourth pair requires one change to the order of the step.
The natural order computes $C$, then $D$, then $DA$, then $CB$, which chains only
$D$ into $DA$.
We compute $D$ and $DA$ first and then $C$ and $CB$, which chains both.
The two orders are equivalent because $C$ and $D$ depend only on $x_3$ and $z_3$ and
feed different multiplications.

The inversion contains far longer chains of a single operation.
We evaluate $z^{p-2}$ with a fixed addition chain whose runs of consecutive
squarings have lengths 1, 2, 1, 5, 10, 20, 10, 50, 100, 50 and 5, separated by 11
multiplications.
Within a run we rename the four output registers that differ from the input set and
invoke the patch again, so every intermediate of the run stays in registers.
Only the first squaring of a run loads a field element and only the last stores one,
which removes 486 of the 508 field-element transfers that the same chain would perform
through a per-operation interface.
We emit the whole inversion as one assembly block, because the register state must
survive from the beginning of a run to its end and a function boundary would not
permit that.

\subsubsection{Patch layout and verification}
\label{sec:c25519-verification}

We validate the installed patches rather than a model of them.
The wrappers declare every register that a patch clobbers, and each wrapper keeps
its state or output pointer in a register that the patch leaves intact, so that a
missing clobber declaration would be observed as a wrong result rather than as a
silent miscompilation.

We test the complete X25519 implementation rather than the two patched operations
in isolation.
The tests comprise the two standalone vectors of RFC~7748, the single-iteration
vector, and the 1000-iteration chain, and they exercise scalar clamping, all 255
ladder steps, the inversion, the final reduction and the encoding.
We run the identical tests for the hand-written C, fiat-crypto, CryptOpt,
\texttt{amd64-51} and \texttt{amd64-64} contenders of \Cref{sec:c25519-perf}, and we
compare the microcode output against each native backend after the chain has
completed.

The 1000-iteration chain also serves a second purpose.
It is our gate on the limb bound discipline described at the end of
\Cref{sec:c25519-mul}.
Each iteration feeds the encoded output of the previous one back as the next scalar
and point, so an output limb that exceeded its permitted range would propagate into
the following scalar multiplication and diverge from the reference within a few
iterations rather than producing a result that happens to agree on one input.
The same chain gates the substitution of the \texttt{amd64-51} assembly field
operations into our ladder, which is only sound if the two implementations agree on
how much headroom a returned limb may consume.

Neither patch contains a conditional operation or an indexed memory access.
Every operand is a register named in the patch image, every shift amount is an
immediate fixed at installation time, and the schedule of each output position is
identical regardless of operand values.
The surrounding ladder performs its conditional swap arithmetically and iterates a
fixed 255 times.
No data-dependent branch or data-dependent address therefore arises, and as in
\Cref{sec:keccak} we assume that the underlying $\mu$ops have data-independent
latency.
